% Part: normal-modal-logic
% Chapter: filtrations
% Section: introduction

\documentclass[../../../include/open-logic-section]{subfiles}

\begin{document}

\olfileid[id]{nml}{fil}{int}

\olsection{Pengantar}

Salah satu pertanyaan penting mengenai suatu logika selalu ialah apakah logika
itu dapat diputuskan, yakni apakah terdapat prosedur efektif yang akan menjawab
pertanyaan ``apakah !!{formula} ini valid?'' Logika proposisional dapat
diputuskan: kita dapat menguji secara efektif apakah !!a{formula} merupakan
tautologi dengan menyusun tabel kebenaran, dan untuk suatu !!{formula} tertentu,
tabel kebenarannya berhingga. Namun, kita tidak dapat begitu saja menguji apakah
suatu formula modal bernilai benar dalam semua model, sebab terdapat tak
berhingga banyak model. Kita dapat mendaftar semua model berhingga yang relevan
bagi suatu formula tertentu, sebab yang relevan hanyalah penetapan himpunan
bagian dunia kepada !!{propositional variable}s yang benar-benar muncul dalam
!!{formula} tersebut. Jika relasi aksesibilitas ditetapkan, berbagai penetapan
$V(p)$ yang mungkin hanyalah semua himpunan bagian dari~$W$, dan jika
$\card{W} = n$, terdapat $2^n$ penetapan semacam itu. Jika !!{formula}~$!A$
memuat $m$ !!{propositional variable}s, terdapat $2^{nm}$ model berbeda
dengan~$n$ dunia. Untuk setiap model itu, kita dapat menguji apakah~$!A$
bernilai benar pada semua dunia cukup dengan menghitung nilai kebenaran~$!A$
pada masing-masing dunia. Tentu saja, kita juga harus memeriksa semua relasi
aksesibilitas yang mungkin, tetapi pada~$n$ dunia pun hanya terdapat berhingga
banyak relasi (tepatnya, sebanyak himpunan bagian dari $W \times W$, yakni
$2^{n^2}$).

Jika yang kita minati bukan logika~\Log{K}, melainkan logika yang didefinisikan
oleh suatu kelas model (misalnya, model refleksif dan transitif), kita juga
harus dapat menguji apakah relasi aksesibilitasnya merupakan jenis yang tepat. Kita dapat
melakukannya apabila kerangka yang kita minati dapat didefinisikan oleh
!!{formula}s modal (misalnya, dengan menguji apakah \Ax{T} dan \Ax{4} valid
pada kerangka itu). Jadi, gagasannya ialah menelusuri semua kerangka berhingga,
menguji setiap kerangka untuk menentukan apakah kerangka itu termasuk dalam
kelas yang kita minati, lalu mendaftar semua model yang mungkin pada kerangka
tersebut dan menguji apakah $!A$ bernilai benar dalam masing-masing model. Jika
tidak, berhentilah: $!A$ tidak valid dalam kelas model yang diminati.

Ada masalah dengan gagasan ini: kita tidak tahu kapan, jika pernah, kita dapat
berhenti mencari. Jika !!{formula} tersebut mempunyai model tandingan
berhingga, prosedur kita akan menemukannya. Namun, jika formula itu tidak
mempunyai model tandingan berhingga, kita tidak akan memperoleh jawaban.
!!{formula} tersebut mungkin valid (sama sekali tidak mempunyai model
tandingan), atau mungkin hanya mempunyai model tandingan tak berhingga yang
tidak akan pernah kita periksa. Masalah ini dapat diatasi jika kita dapat
menunjukkan bahwa setiap !!{formula} yang mempunyai model tandingan juga
mempunyai model tandingan berhingga. Jika demikian, kita mengatakan bahwa
logika tersebut mempunyai \emph{sifat model berhingga}.

Namun, bagaimana kita menunjukkan bahwa suatu logika mempunyai sifat model
berhingga? Salah satu caranya ialah menemukan cara untuk mengubah suatu model
(tandingan) tak berhingga bagi~$!A$ menjadi model berhingga. Jika hal itu dapat
dilakukan, setiap kali terdapat model tempat $!A$ tidak bernilai benar, model
berhingga yang dihasilkan juga membuat $!A$ tidak bernilai benar. Model
berhingga itu akan muncul dalam daftar semua model berhingga kita, dan pada
akhirnya kita akan menentukan, bagi setiap formula yang tidak valid, bahwa
formula itu memang tidak valid. Prosedur kita tidak akan berhenti jika formula
tersebut valid. Jika selain itu kita dapat menunjukkan bahwa terdapat suatu
ukuran maksimum bagi model berhingga yang dihasilkan prosedur kita, dan bahwa
ukuran maksimum tersebut hanya bergantung pada !!{formula}~$!A$, kita akan
mempunyai batas ukuran model berhingga yang perlu diuji dalam pencarian model
tandingan. Jika sampai batas itu kita belum menemukan model tandingan, berarti
tidak ada model tandingan. Dengan demikian, prosedur kita benar-benar akan
memutuskan pertanyaan ``apakah $!A$ valid?'' bagi setiap !!{formula}~$!A$.

Strategi yang sering berhasil untuk mengubah struktur tak berhingga menjadi
struktur berhingga ialah ``mengidentifikasi'' !!{element}s struktur yang
berperilaku sama dalam segi-segi yang relevan. Jika terdapat tak berhingga
banyak dunia dalam~$\mModel{M}$ yang berperilaku sama dalam segi-segi yang
relevan, kita mungkin dapat mengumpulkan \emph{semua} dunia ke dalam berhingga
banyak ``kelas'' dunia semacam itu (yang mungkin saja tak berhingga). Dengan
kata lain, kita harus mempartisi himpunan dunia dengan cara yang tepat, yakni
sedemikian rupa sehingga setiap blok partisi mungkin memuat tak berhingga
banyak dunia, tetapi hanya terdapat berhingga banyak blok. Kemudian kita
mendefinisikan model baru~$\mModel{M^*}$ yang dunia-dunianya adalah blok-blok
partisi tersebut. Berhingga banyak blok partisi dalam model lama memberi kita
berhingga banyak dunia dalam model baru, yakni suatu model berhingga. Mari
sebut blok partisi yang memuat dunia~$w$ sebagai $[w]$. Kita menghendaki agar
$\mSat{M}{!A}[w]$ jika dan hanya jika $\mSat{M^*}{!A}[{[w]}]$, sebab kita
menghendaki model baru menjadi model tandingan bagi~$!A$ jika model lama
demikian. Hal ini menuntut kita mendefinisikan partisi dan relasi aksesibilitas
dari~$\mModel{M^*}$ dengan cara yang tepat.

Untuk melihat bagaimana cara ini berjalan, pertama-tama bayangkan bahwa relasi
aksesibilitasnya universal. $\mSat{M}{\Box !B}[w]$ jika dan
hanya jika untuk setiap $v \in W$, $\mSat{M}{!B}[v]$, dan hal yang sama
berlaku bagi $\mModel{M^*}$, hanya saja dengan $[w]$ dan $[v]$. Sebagai
gagasan awal, katakanlah bahwa dua dunia $u$ dan $v$ ekuivalen (termasuk dalam
blok partisi yang sama) jika keduanya bersepakat mengenai semua
!!{propositional variable}s dalam~$\mModel{M}$, yakni $\mSat{M}{p}[u]$ jika
dan hanya jika $\mSat{M}{p}[v]$. Misalkan
$V^*(p) = \Setabs{[w]}{\mSat{M}{p}[w]}$. Tujuan kita ialah menunjukkan bahwa
$\mSat{M}{!A}[w]$ jika dan hanya jika $\mSat{M^*}{!A}[{[w]}]$. Jelas bahwa
kita akan membuktikannya melalui induksi: kasus dasarnya ialah $!A \equiv p$.
Pertama-tama, andaikan $\mSat{M}{p}[w]$. Maka $[w] \in V^*(p)$ menurut
definisi, sehingga $\mSat{M^*}{p}[{[w]}]$. Sekarang andaikan
$\mSat{M^*}{p}[{[w]}]$. Ini berarti bahwa $[w] \in V^*(p)$, yakni bagi suatu
$v$ yang ekuivalen dengan~$w$, $\mSat{M}{p}[v]$. Namun, ``$w$ ekuivalen
dengan $v$'' berarti ``$w$ dan $v$ membuat !!{propositional variable}s yang
sama bernilai benar,'' sehingga $\mSat{M}{p}[w]$. Sekarang, untuk langkah
induktif, misalnya $!A \ident \lnot !B$. Maka
$\mSat{M}{\lnot !B}[w]$ jika dan hanya jika $\mSat/{M}{!B}[w]$, jika dan
hanya jika $\mSat/{M^*}{!B}[{[w]}]$ (menurut hipotesis induksi), jika dan
hanya jika $\mSat{M^*}{\lnot !B}[{[w]}]$. Hal yang sama berlaku bagi operator
nonmodal lainnya. Cara ini juga berlaku bagi $\Box$: andaikan
$\mSat{M^*}{\Box !B}[{[w]}]$. Ini berarti bahwa bagi setiap $[u]$,
$\mSat{M^*}{!B}[{[u]}]$. Menurut hipotesis induksi, bagi setiap $u$,
$\mSat{M}{!B}[u]$. Akibatnya, $\mSat{M}{\Box !B}[w]$.

Dalam kasus umum, ketika kita juga harus mendefinisikan relasi aksesibilitas
bagi $\mModel{M^*}$, keadaan menjadi lebih rumit. Kita akan menyebut suatu
model~$\mModel{M^*}$ sebagai \emph{filtrasi} jika relasi
aksesibilitasnya~$R^*$ memenuhi syarat-syarat yang diperlukan agar pembuktian
induktif di atas berjalan. Dengan demikian, setiap filtrasi~$\mModel{M^*}$
akan membuat $!A$ bernilai benar pada $[w]$ jika dan hanya jika
$\mModel{M}$ membuat $!A$ bernilai benar pada~$w$. Namun, sekarang kita juga
harus menunjukkan bahwa filtrasi \emph{memang ada}, yakni bahwa kita dapat
mendefinisikan~$R^*$ sehingga memenuhi syarat-syarat yang diperlukan. Agar
hal ini berhasil, kita harus mensyaratkan agar dunia $u$ dan~$v$ dianggap
ekuivalen bukan hanya ketika keduanya bersepakat mengenai semua
!!{propositional variable}s, melainkan mengenai semua sub-!!{formula}s
dari~$!A$. Karena $!A$ hanya mempunyai berhingga banyak sub-!!{formula}s,
hal ini tetap menjamin bahwa filtrasi tersebut berhingga. Ada lebih dari satu
cara untuk mendefinisikan filtrasi, dan untuk memastikan bahwa relasi
aksesibilitas filtrasi mempunyai sifat-sifat yang diperlukan (misalnya
refleksif, transitif, dan sebagainya), kita harus kreatif dalam mendefinisikan
$R^*$.


\end{document}
